\documentclass[a4paper, 10.5pt, oneside]{report}
%------------PREAMBLE
\input{preamble}
\usepackage{morefloats}
%---------------------------
                                          
\begin{document}

%---------------------------TITLE PAGE
\begin{titlepage}
\begin{center}

% Upper part of the page. The '~' is needed because \\
% only works if a paragraph has started.
\includegraphics[width=0.35\textwidth]{figures/uom_banner.pdf}~\\[0.3cm]

\textsc{\Large Final Report}\\[0.3cm]

% Title
\HRule \\[0.3cm]
{ \huge \bfseries The Art of Scientific Computing:\\
Project Title: A selection of granary queueing models }\\[0.3cm]

\HRule \\[0.8cm]

% Author and supervisor
\begin{minipage}{0.45\textwidth}
\begin{flushleft} \large
\textit{Author:}\\
William Batchelor
\end{flushleft}
\end{minipage}
\begin{minipage}{0.45\textwidth}
\begin{flushright} \large
\textit{Subject Handbook code:}\\
COMP-90072
\end{flushright}
\end{minipage}
\\[1.0cm]

\includegraphics[width=0.50\textwidth]{figures/coulomb_direct.pdf}~\\[0.5cm]

\vspace{1.0cm}

% Place
Faculty of Science\\
The University of Melbourne
\\[0.5cm]

% Submission date
{\small Original Version: 1st March 2018}
\\[1.5cm]

% Lecturer
Subject co-ordinators: A/Prof. Suzie Sheehy \& Prof. Christian Reichardt
\\[0.8cm]


% Bottom of the page
\end{center}
\end{titlepage}

%---------------------------

%---------------------------TABLE OF CONTENTS
\tableofcontents
%---------------------------

%---------------------------MAIN BODY

\chapter{Introduction}
Granaries are storage structures for grain, used to protect harvested grain from poor weather and pests.
To prevent spoilage grain needs to be kept at a low moisture content; this state may initially be achieved
by harvesting during dry conditions, or by drying harvested grain. Drying sometimes happens at the granary.

It is typical for a granary to service all of the farms in an area, often operating either as a collective or
under commercial agreements.

Grain must be harvested at the correct time of year, and the need for dry conditions leads to busy peak periods
for granaries corresponding with appropriate weather. Deliveries cluster, as farmers restart harvesting
operations after inclement weather or race to harvest as much as possible before rain.

These high demand periods can lead to significant queues, incurring costs of time and money as trucks sit idle
and farms run out of buffer storage. Understanding the likely demands placed on a granary is therefore important
to decisions regarding its design and operation. Planners may wish to invest in equipment to reduce unloading
time or add additional service bays; managers may wish to alter policy regarding opening hours or the order in
which trucks are served.

Mathematical models are a useful way of investigating these demands and evaluating different investments and
policies. In this investigation we implement three simulation models of a granary, and investigate where the models
agree and disagree. We compare to closed form results for queueing systems, and consider where the models likely
fall short of reality.



\chapter{Background}
In which we cover basic queueing theory, classify our problem and mention
the closed form results that are going to be important for verification.

Before considering the specifics of granaries, it is worth establishing some basic queueing theory.

In the typical terminology, a customer is an entity subject to a queue, and a server is an entity that
provides some resource to customers. In the usual case a queue is FIFO - First In First Out, a system familiar
to anyone who has visited a grocery store. It is worth being aware of an alternative structure, the priority queue,
where customers have a priority value, with servers always serving the next highest priority customer.

Queueing systems are classified with Kendall's notation, a shorthand to quickly communicate the behaviour of the system.
A classification is shown as A/S/c, where A denotes the probability distribution driving customer arrival,
S denotes the distribution driving service time, and c denotes the number of servers in the system.

We will consider some examples to get an idea of the scheme.

A D/D/1 system has deterministic arrivals, deterministic service time, and a single server. In other words, customers
always show up the same amount of time after the previous customer, and they are always served in a fixed amount
of time. This system is of little interest, if the server processes customers faster than they arrive
there is never a queue, and if it is slower the queue grows without bound.

An M/D/1 system has Poisson arrivals, deterministic service times, and a single server. M stands for memoryless; Poisson
arrivals means the time between arrivals is an exponential random variable.

An M/M/1 system possesses Poisson arrivals, exponentially distributed service times, and a single server.

An M/M/c system possesses Poisson arrivals, exponentially distributed service times, and c servers.

The nice thing about all of the above systems is they possess simple service functions and memoryless (Poisson) arrival functions,
permitting simple closed form solutions to the behaviour of the queue in the long term.

A G/D/1 queue on the other hand, with a general (arbitrary) distribution for arrivals, does not necessarily
have a closed form solution, meaning analysis must rely on approximation or simulation. For completeness it is worth noting that
M/G/1 queues do possess closed form results; memoryless arrival is the important condition for closed form solutions to exist.

Need words on closed form for M/D/1 - Pollaczek-Khinchine

Finite source and extended Kendall needs some words

Little's law






\chapter{Models}
In which we discuss the three models of a granary we're going to be using

First the shared infrastructure of a simple queue and a single deterministic server

Before discussing the three specific models, we will touch on their commonalities, and the behaviour
we want to capture. All models feature a single FIFO queue at the granary determining service order.
Further, service time is deterministic.

Clustering due to weather
The weather system


\section{Open model}
The first model is the simplest. At its simplest it is an M/D/1 queueing system, with Poisson arrivals.

Interarrival time is driven by a parameter $\lambda$, trucks arrive, queue, unload at the granary, before
disappearing into the aether.

We can run the simulation in this simple form to confirm against closed form results, but in practice
we want to incorporate opening hours and weather effects.

To that end, we track the moisture level in the grain throughout the season. Any trucks arriving while
moisture is above some threshold value are rejected with linearly increasing probability. At a certain
point, all arrivals are rejected. This has the effect of increasing the interarrival time as conditions
become unfavourable to the harvest. Similarly, if the granary is closed for the night,
all arrivals are rejected.


\section{Closed model with abstract farm}
The second model increases fidelity by closing the loop and no longer consigning trucks to the void.
Instead, we introduce a new entity with its own queue, the farm. Grain accumulates at the farm as
it is harvested, trucks queue there to pick up grain, travel to the granary to unload it, before
returning to the farm to queue again.

There is a fixed fleet of trucks available, each truck possesses its own travel time between the granary
and the farm. Travel times are drawn uniformly from a circle, accounting for the fact that area of farmland
available increases as the square of the distance from the granary. Each trip noise is applied to the travel
time reflecting variance in conditions on the road.

Again we must deal with weather and opening hours.

Weather enters this model through the rate at which grain is harvested. Similarly to the open model, below
a moisture threshold harvest proceeds at full speed, above it the rate decreases linearly until it eventually
reaches zero.

Opening hours enter this model in the step where trucks attempt to depart the farm and head to the granary.
The decision was made that farms would operate continuously except as prevented by moisture levels, whereas
the granary keeps to its opening hours, except for serving trucks in its queue at closing time,
which we allow as overtime. To that end, trucks expected to arrive after closing time are instead scheduled
to arrive the next morning. This matches morning clustering often seen in reality.

We also wish to promote clustering after a stoppage due to rainfall. To that end, when such an event occurs,
trucks are prohibited from leaving the granary until a threshold of stored grain is reached. The intention
here is to reflect that this farm is intended as an abstraction of a number of farms. If all of the farms
possess between them only enough grain to fill a single truck, no truck will be filled. Assuming they harvest
at a similar rate, all of them would be expected to release trucks at similar times; thus, the threshold. Away from
stoppages it's assumed that if the abstract farm has grain to dispatch then the hypothetical farm any given
truck is serving would have grain to dispatch after its round trip. In the case of trucks with short round trips
that may not be a good assumption.

Another flaw of this model is the bottleneck around truck dispatch. If the time it takes to fill an individual truck
is set to a plausible level, then with a single server the rate at which grain can be dispatched is limited compared
to a model with multiple farms, each with their own server, 






\section{Closed model with farms}
Our third model partakes of most of the characteristics of the second model. The crucial difference is this
time we consider multiple farms with their own fleet of trucks and their own queue; the other differences
all follow from this.

Travel time to and from each farm is drawn uniformly from a circle, each truck in a farms fleet possesses the
same travel time. Noise is still applied independently to each trip.

Weather is treated the same, except each farm of course has its own grain. W`e no longer need to enforce
a threshold of grain before releasing trucks after a stoppage, each farm will naturally wait until it has
sufficient grain to fill a truck before dispatching one, producing the desired clustering.

The key difference between 


\chapter{Implementation}
In which we discuss the implementation in python. Discuss the initial simpy profiling,
and the subsequent handmade simulation with a queue using a deque and a priority
queue of upcoming events using a heapq. Discuss parallelisation with joblib - embarassingly parallel


\chapter{Validation}
Out come the closed form results, most useful for model 1. In correct regimes
model 2 and 3 should converge to something like those results as well


\chapter{Experiments}
How do mean and tail waiting times change as we increase weather variance?

How does overtime increase as we increase weather variance?

How do the models compare to one another?

Report on parallelisation efficiency


\chapter{Results and Discussion}
When models are tuned to be similar, do the tunings make sense analytically?

Where do models disagree when tuned to be similar?


\chapter{Conclusion}





\end{document}